\section{Introduction}
\label{sec:intro}

This document describes the physics, numerical methods, and modeling assumptions
behind the two central calculations in the WOTT cycling performance application:

\begin{enumerate}
    \item \textbf{Forward race simulation}, implemented by \nolinkurl{IPCalculator}:
    given a rider's physical parameters, a target power plan, and environmental and
    track conditions, integrate the equations of motion forward in time to predict
    velocity, position, split times, and power output over the course of a race.

    \item \textbf{Aerodynamic drag fitting}, implemented by \nolinkurl{CdAFitter}:
    given recorded speed and power data from a \nolinkurl{.fit} file (the standard
    recording format written by GPS/ANT+ cycling computers), invert the same
    equations of motion to estimate a rider's aerodynamic drag area, \CdA{}, from
    field data.
\end{enumerate}

Both calculations share a single underlying physical model: a point mass moving
along a one-dimensional track coordinate $s$, subject to rolling resistance,
aerodynamic drag, and a rider-supplied propulsive force, together with an optional
quasi-static cornering correction that activates whenever track curvature is
present, with further geometric and energetic refinements activating when the
rider's center-of-mass height is also known. \Cref{sec:simulation-model} develops this
equation of motion and its numerical solution. \Cref{sec:track-geometry}
describes how the track's curvature profile is constructed from a small set of
geometric parameters. \Cref{sec:cornering-model} then derives the cornering
correction in detail, since it is the most involved part of the model and
depends on that curvature profile. \Cref{sec:cda-fitting}
shows how the forward model is inverted to estimate \CdA{} from measured ride
data. \Cref{sec:assumptions} collects the simplifying assumptions made throughout
and states their limitations plainly. \Cref{sec:notation} tabulates every symbol
used in this document, its SI unit, its default value where one exists, and the
corresponding function or variable name in the implementation
(\nolinkurl{calcs.py}).

\subsection*{State and notation conventions}

The forward simulation's state at time $t$ is the pair
\begin{equation}
    y(t) = \bigl(v(t),\ s(t)\bigr)\ ,
    \label{eq:state-vector}
\end{equation}
where $v$ is the rider's speed along the direction of travel, in
\si{\metre\per\second}, and $s$ is the distance traveled along the track, in
\si{\metre}. All quantities in this document are expressed in SI units unless
stated otherwise.
